Program analysis has a long history in analyzing and improving the program's execution behavior.
Since few decades ago, when computer programs were only used for research or business purpose and the expectation of computer programs was just to execute correctly without bugs,
the most critical behaviors with respect to programs is the execution correctness behavior.
To this sense, the program analysis community pays much attention on analyzing functional correctness properties.
It determines whether a program executes correctly without getting stuck because of bugs.
% The execution correctness can be proved by showing the functional correctness property of the program with the help of some formal verification methods such as type system and program logic.
% 
However, the expectation has also changed along with the popularity of mobile devices and wide applications of big data. 
It is not enough for these programs running on mobile phones or those programs handling big data, 
to just run without bugs. 
The non-functional execution properties come into play in the new era.
We think programs on mobile devices are of great significance in today's life.
To provide people with high-quality mobile devices
service through programs,
lots of programs' non-functional properties are worth to be studied. This thesis studies two non-functional properties in the \emph{quantitative-reachability} property category which can significantly influence people's daily life in terms of the \emph{convenience},
\emph{enjoyment} and the \emph{efficiency}.


\emph{Convenience And Enjoyment.} We notice that people's daily life
quality is closely related to the experience of using the applications on their mobile devices.
Numerous mobile device applications
such as video or music players, online games, shopping, etc.,
provide users with personalized recommendations, to provide user with a more convenient and enjoyable experience.
In this sense, part of the convenience and enjoyment in mobile experience
is influenced by
the recommendation accuracy when using applications.
The accuracy of the recommendation
relies heavily on the
outputs of the programs which implement some
data analysis over big data.
This brings my attention to
some non-functional properties
that can help to improve the output accuracy of
the programs that implement a data analysis.
We look at the scenario
where a data analysis program
analyzes a sample data
and computes some generalized results
on the unknown population data
where the sample is drawn.
Then, these generalized results are used to provide the
users who come from the unknown population with personalized recommendations.
Some users occasionally receive inaccurate recommendations,
because they are not in the sample data but get services by using the
same
%  generalized 
analysis results.
In these situations, the generalization error comes into play. 
It measures the difference between the 
data analysis result on a sample data 
and the ideal analysis result expected to be
achieved by
analyzing the entire population.
This error shows the reliability of the data analysis and reflecting the accuracy of the personalized recommendations.
High generalization error makes the analysis results
not reliable and reduces the accuracy on a large scale.
Studies found that some non-functional properties of data analysis algorithms can help to control the generalization error, especially when the data analysis is adaptive.
\\
These non-functional properties can be expressed from two perspectives, the \emph{quantitative} and \emph{reachability}.
We are interested in analyzing them as \emph{quantitative-reachability} properties using the principles of programming languages to help control the generalization error of corresponding algorithm.
In this sense, we can help to improve accuracy for the recommendation results computed by these algorithms
and guarantee the \emph{convenience} and the \emph{enjoyment} of people's life.
We specifically look into the \emph{adaptivity rounds} property in our \redd{PART \romannum{1}}. It captures the generalization error in terms of whether the queries are dependent on each other (the \emph{reachability}) and how many times they are dependent (the \emph{quantitative}).
Under the light of existing machine learning studies, we want to use this \emph{quantitative-reachability} property to control the generalization error.


\emph{Efficiency.} High-quality modern daily life
is not only
determined by the enjoyment or convenience in using some mobile device applications,
but also
the performance, 
privacy, etc., of the applications on a mobile device or software on a computer. 
Suppose we are playing an online game on our smartphones, what do we care about? 
In games such as racing, 
we care about the performance of the game, in another word, if it runs fast, if it crashes due to being out of memory, and so on.
In some other strategy games, we also care about whether my strategy is leaked to
other parties, i.e., the privacy quality of the game.
We observe that the performance of
a program is usually related to resource cost, which is
mainly determined by
whether some lines of code are executed and how times they
are executed.
From the same perspective, we observe that whether the data is leaked is also related to
whether certain pieces of the program code are executed,
and their execution times.
This brings my attention to other non-functional properties,
the
reachability of certain lines of program code, and the times that certain code is executed.
\\
These two properties can indeed be captured by a \emph{quantitative-reachability} property from two perspectives, whether certain piece of code is executed (the \emph{reachability}) and how many times it is executed  (the \emph{quantitative}).
In \cite{GulwaniJK09}, this property is referred to as \emph{reachability-bound}, which is a worst-case bound on the number of times a given control location 
inside a procedure is visited during the program execution.
It has broad applications in analyzing other program property beyond itself.
For example in the resource cost analysis, this \emph{quantitative-reachability} property
can help to provide a tighter
bound on the resources consumed by a program such as time, memory,
network traffic, power, etc.
In the data analysis area, for example, when analyzing the program's \emph{adaptivity}
as introduced above,
the \emph{reachability-bound} on each program control location
can improve the precision of the estimated bound.
Analyzing and controlling it from two perspectives can better help to improve program reliability, security, privacy, usability, etc., and guarantee the high-quality
service through the mobile device people's life than just from single view. 
Motivated by this, also as a continuation in improving the precision when analyzing the \emph{adaptivity rounds} in the first setting, we look into the \emph{reachability-bound} in our \redd{PART \romannum{2}}


For both of the two conceptual \emph{quantitative-reachability} properties, we first formalize them  through specific programming language models.
Then we develop new static program analysis frameworks to best exploit the benefit of these properties in data analysis and resource usage. 
However, to fully utilize the aforementioned benefits of the new analysis frameworks on these non-functional properties, the appropriate implementation is inevitable. 
To this end, we also take one step in algorithmitizing and implementing the analysis frameworks
which is used to formalize and statically estimate these properties.
The evaluation results over real world benchmarks report the effectiveness of our work in real world and its potential applications beyond the program analysis area.

% Last but not least, reasoning these \emph{quantitative-reachability} properties are not limited to static program analysis. 
% We are not only interested in programs implementing the specific algorithm, but also willing to study the properties of the algorithm itself.
